%% ma: L12-B04-0007
%% lop: 12
%% chuong: 1
%% bai: 4
%% dang: 12.4.7
%% loai: TLN
%% muc: VD
%% dap_an: "19"
%% nguon: "Ảnh giáo viên gửi 11/10/2026 (Đề kiểm tra giữa HK1 Toán 12, THPT Việt Đức 2026-2027, mã đề 0103), Phần III Câu 5"
%% kien_thuc: "Bài 4 (đồ thị hàm phân thức bậc nhất trên bậc nhất: tiệm cận, giao điểm với các trục), Bài 3 (tiệm cận)"
%% trang_thai: chinh-thuc
%% ghi_chu: "Đã duyệt 12/10/2026."
\begin{ex}
Cho hàm số $y=\dfrac{3x+b}{cx+d}$ với $b,c,d$ là hằng số, có đồ thị như hình vẽ.
\begin{center}
\begin{tikzpicture}[x=.7cm,y=.55cm]
  \clip (-6.6,-1.8) rectangle (7.6,9.2);
  \draw[dashed,gray] (-2,-1.8)--(-2,9.2);
  \draw[dashed,gray] (-6.6,3)--(7.6,3);
  \draw[thick,domain=-6.3:-2.2,samples=80] plot (\x,{(3*\x+12)/(\x+2)});
  \draw[thick,domain=-1.7:7.4,samples=80] plot (\x,{(3*\x+12)/(\x+2)});
  \draw[truc] (-6.6,0)--(7.6,0) node[below]{$x$};
  \draw[truc] (0,-1.8)--(0,9.2) node[right]{$y$};
  \node[below,font=\footnotesize] at (-4,0) {$-4$};
  \node[below,font=\footnotesize] at (-1.55,0) {$-2$};
  \node[left,font=\footnotesize] at (0,3.3) {$3$};
  \node[left,font=\footnotesize] at (0,6.3) {$6$};
  \node[below left,font=\footnotesize] at (0,0) {$O$};
  \node[font=\footnotesize] at (4.2,4.7) {$y=f(x)$};
\end{tikzpicture}
\end{center}
Biểu thức $T=2b+3c-4d$ có giá trị bằng bao nhiêu?
\shortans{19}
\loigiai{Đồ thị có tiệm cận ngang $y=3$ nên $\dfrac3c=3$, suy ra $c=1$. Tiệm cận đứng $x=-2$ nên $-2c+d=0$, suy ra $d=2$. Đồ thị cắt trục hoành tại $(-4;0)$ nên $3\cdot(-4)+b=0$, suy ra $b=12$ (kiểm tra: giao điểm với $Oy$ là $\left(0;\dfrac bd\right)=(0;6)$ khớp hình). Vậy $T=2\cdot12+3\cdot1-4\cdot2=19$.}
\end{ex}
